%!TEX root = ../main.tex

\addchap{Anhang}
\begin{itemize}
	\item \ref{apx:lighthouse} Lighthouse-Messergebnisse \dotfill{} \pageref{apx:lighthouse}
\end{itemize}

\chapter{Lighthouse-Messergebnisse}
\label{apx:lighthouse}

Die Performance-Analyse wurde mit Google Lighthouse Version~12 auf dem produktiven Firebase-Hosting-Endpunkt nach einem vollständigen Produktionsbuild durchgeführt. Alle Messungen erfolgten im Inkognito-Modus von Google Chrome (Version 131) ohne aktive Browser-Erweiterungen, um Verfälschungen durch lokale Cache-Inhalte oder Drittanbieter-Skripte auszuschließen. Die Netzwerkkonditionierung entsprach der Lighthouse-Voreinstellung \textit{Mobile -- Simulated Slow 4G}.

\begin{table}[H]
\centering
\begin{tabularx}{\textwidth}{|X|c|c|c|}
\hline
\textbf{Metrik} & \textbf{Zielwert} & \textbf{Startseite} & \textbf{Produktdetailseite} \\
\hline
Performance Score & $\geq$ 90 & 91 & 88 \\
\hline
Largest Contentful Paint (LCP) & $\leq$ 2{,}5\,s & 1{,}8\,s & 2{,}3\,s \\
\hline
Interaction to Next Paint (INP) & $\leq$ 200\,ms & 80\,ms & 95\,ms \\
\hline
Cumulative Layout Shift (CLS) & $\leq$ 0{,}1 & 0{,}02 & 0{,}04 \\
\hline
Best Practices Score & 100 & 100 & 100 \\
\hline
\end{tabularx}
\caption{Lighthouse-Messwerte der produktiven Anwendung}
\label{tab:lighthouse-ergebnisse}
\end{table}

Der geringfügig niedrigere Performance-Score der Produktdetailseite (88 gegenüber 91) ist auf die erhöhte LCP-Zeit zurückzuführen, die durch den asynchronen Firestore-Datenabruf entsteht: Der Hauptinhalt der Seite kann erst nach Abschluss des Datenbankaufrufs gerendert werden, was eine strukturell bedingte Latenz einführt. Eine serverseitige Vorrendering-Lösung (z.\,B. Nuxt.js mit SSR) würde diesen Wert verbessern, überschreitet jedoch den Rahmen der vorliegenden Studienarbeit.
